\section*{अध्याय \ref{ch:g3:stones-pebbles-sand} --- पत्थर, कंकड़ और रेत}

\begin{solution}{exo:g3:stones-pebbles-sand:1}
\omterm{def:g3:stones-pebbles-sand:rock}{आधार चट्टान} वह ठोस \omterm{def:g3:stones-pebbles-sand:rock}{चट्टान} है जो हर जगह मिट्टी के नीचे फैली है। हम उसे
वहाँ देखते हैं जहाँ उसे छिपाने के लिए मिट्टी नहीं है: \omterm{def:g1:landscapes-around-us:landform}{पर्वतों} की खड़ी
ढलानों पर, \omterm{def:g2:water-on-earth:ocean}{समुद्र} किनारे की खड़ी \omterm{def:g3:stones-pebbles-sand:rock}{चट्टानों} में, \omterm{def:g1:landscapes-around-us:landform}{पहाड़ी} को काटकर बनी सड़कों के
किनारे।
\end{solution}

\begin{solution}{exo:g3:stones-pebbles-sand:2}
\omterm{def:g3:stones-pebbles-sand:sand}{कीचड़}, \omterm{def:g3:stones-pebbles-sand:sand}{रेत}, \omterm{def:g3:stones-pebbles-sand:pebble}{कंकड़}, शिलाखंड।
\end{solution}

\begin{solution}{exo:g3:stones-pebbles-sand:3}
1 मिलीमीटर 2 मिलीमीटर से छोटा है: \omterm{def:g3:stones-pebbles-sand:sand}{रेत}। 3 मिलीमीटर 2 मिलीमीटर से बड़ा है:
\omterm{def:g3:stones-pebbles-sand:pebble}{बजरी}।
\end{solution}

\begin{solution}{exo:g3:stones-pebbles-sand:4}
प्राकृतिक: \omterm{def:g3:stones-pebbles-sand:pebble}{कंकड़}, ग्रेनाइट का ब्लॉक (लोगों ने काटा है, पर \omterm{def:g3:stones-pebbles-sand:rock}{चट्टान} प्राकृतिक
है), \omterm{def:g2:water-on-earth:ocean}{समुद्र} \omterm{def:g1:landscapes-around-us:water}{तट} की \omterm{def:g3:stones-pebbles-sand:sand}{रेत}। लोगों के बनाए: ईंट और काँच।
\end{solution}

\begin{solution}{exo:g3:stones-pebbles-sand:5}
\omterm{def:g1:landscapes-around-us:water}{नदी} के पानी ने उसे बहुत समय तक तली पर लुढ़काया है और दूसरे पत्थरों से
टकराया है: उसके कोने घिस गए हैं।
\end{solution}

\begin{solution}{exo:g3:stones-pebbles-sand:6}
ग्रेनाइट, दीवार के लिए; पत्थर की पतली, चपटी पटियाँ (स्लेट जैसी \omterm{def:g3:stones-pebbles-sand:rock}{चट्टान}), छतों के लिए।
\end{solution}

\begin{solution}{exo:g3:stones-pebbles-sand:7}
गोल वाले ने: पानी किसी पत्थर को जितना ज़्यादा लुढ़काता है, उसके कोने उतने ही
ज़्यादा घिसते हैं। नुकीला पत्थर उस \omterm{def:g3:stones-pebbles-sand:rock}{चट्टान} से ज़्यादा दूर नहीं गया जिससे वह टूटा था।
\end{solution}

\begin{solution}{exo:g3:stones-pebbles-sand:8}
सबसे बड़ा \omterm{def:g3:stones-pebbles-sand:pebble}{कंकड़} 3 सेंटीमीटर लंबा है। एक सेंटीमीटर 10 मिलीमीटर का होता है,
और $10 \div 2 = 5$: 2 मिलीमीटर के पाँच कण 1 सेंटीमीटर में समा जाते हैं।
\end{solution}

\begin{solution}{exo:g3:stones-pebbles-sand:9}
शांत तालाब में। तेज़ पानी \omterm{def:g3:stones-pebbles-sand:sand}{कीचड़} के नन्हे कणों को चलाता रहता है; वे सिर्फ़
वहीं तली में बैठ पाते हैं जहाँ पानी लगभग ठहरा हुआ हो।
\end{solution}

\begin{solution}{exo:g3:stones-pebbles-sand:10}
क्योंकि वे ज़मीन से लिए गए पदार्थों से बनते हैं: ईंटें मिट्टी से, कंक्रीट
\omterm{def:g3:stones-pebbles-sand:sand}{रेत}, \omterm{def:g3:stones-pebbles-sand:pebble}{बजरी} और सीमेंट से, काँच \omterm{def:g3:stones-pebbles-sand:sand}{रेत} से। मिट्टी, \omterm{def:g3:stones-pebbles-sand:sand}{रेत} और \omterm{def:g3:stones-pebbles-sand:pebble}{बजरी} टूटी \omterm{def:g3:stones-pebbles-sand:rock}{चट्टान} से
आती हैं।
\end{solution}
